\chapter{修改}\label{modifications}

\section{合并后的引擎}

\subsection{变革的必要性}

\topicindex {engines}
\topicindex {history}

\LUATEX\ 的第一个版本仅增加了少量原语，在很大程度上与 \PDFTEX\ 相同。随后我们将 \ALEPH\ 的大量代码合并进来，得到了更多原语。当代码趋于稳定后，我们决定清理该程序相当混杂的本质。这意味着一些原语被提升为核心原语（通常使用不同的名称），另一些则被移除。这使得我们能够开始清理代码库。在 \ref{enhancements} 中我们讨论了一些新原语，这里将介绍大多数经过调整的原语。

除了由新功能引起的预期变化外，还有一些不太预期的变化。这些有时是新（冲突）特性的副作用，或者更常见的是清理内部接口所必需的变更。这些也会一并提及。

\subsection{与 \TEX\ 3.1415926 相比的变化}

\topicindex {\TEX}

当然，一切都始于传统的 \TEX。即使我们从 \PDFTEX\ 起步，大部分仍源自原始版本。但我们有所偏离。

\begin{itemize}\item 当前代码库使用 \CCODE\ 而非 \PASCAL\ 编写。我们尽可能使用 \CWEB。因此，不再是一个大文件加上变更文件，而是按类别组织的多个文件，如 \type {tex}、\type {pdf}、\type {lang}、\type {font}、\type {lua} 等。转换为 \CCODE\ 过程中留下了一些痕迹，但我们会在适当的时候清理源代码并确保文档完善。许多文件采用 \CWEB\ 格式，但其他文件（如与 \LUA\ 接口的文件）则是 \CCODE\ 文件。当然，我们希望尽可能贴近原始版本，以便 Don Knuth 所著的 \TEX\ 基础原理文档仍然适用。

\item 有关段落构建、语言处理和断词方面的许多小变化，请参见 \ref{languages}。最重要的变化是：在单词中间插入花括号组（如 \type {of{}fice}）不再阻止连字的创建。

\item 不再有池文件，所有字符串都在编译期间嵌入。

\item 说明符 \type {plus 1 fillll} 不会产生错误。多余的 ``l 只会被简单排版出来。

\item \prm {endlinechar} 和 \prm {newlinechar} 的上限为 127。

\item 缩放（\prm {mag}）仅在 \DVI\ 输出模式下受支持。你可以设置此参数，它甚至能与 \type {true} 单位一起工作，直到切换到 \PDF\ 输出模式。使用 \PDF\ 输出时，最好不要改动 \prm {mag} 变量。这种模糊行为与最终需要 \DVI\ 输出时却使用 \PDF\ 后端相关功能的情况并无太大区别。

    输出模式被冻结后（通常在第一页 shipped out 时发生）或启用 \PDF\ 输出后，\type {true} 说明将被忽略。当预加载适配 \LUATEX\ 的 plain 格式时，\prm {mag} 参数可能已经被设置。

\item 当 \type {\globaldefs} 为正值而请求局部赋值时，如果 \type {\tracingcommands} 大于 1，日志中会显示 \type {{\global enforced}}。当 \type {\globaldefs} 为负值而通过 \type {\global}、\type {\gdef} 等请求全局赋值时，日志会记录 \type {{\global canceled}}。

\end{itemize}

\subsection{与 \ETEX\ 2.2 相比的变化}

\topicindex {\ETEX}

作为事实上的标准扩展，我们当然提供 \ETEX\ 功能，但做了一些小的调整。

\begin{itemize}\item \ETEX\ 功能始终存在并启用，因此不再需要为 initex\ 添加前置星号或 \type {-etex} 开关。

\item 不包含 \TEXXET\ 扩展，因此原语 \type {\TeXXeTstate}、\type {\beginR}、\type {\beginL}、\type {\endR} 和 \type {\endL} 均不存在。相反，我们以 \OMEGA/\ALEPH\ 的方向性处理方法为起点。

\item \ETEX\ 的 \prm {tracingassigns} 和 \prm {tracingrestores} 输出的部分跟踪信息不存在。

\item \LUATEX\ 中的寄存器管理采用 \OMEGA/\ALEPH\ 模型，因此最大值为 65535，实现使用扁平数组，而非 \ETEX\ 的扁平与稀疏混合模型。

\item 当使用 kpathsea 查找文件时，\LUATEX\ 使用 \type {ofm} 文件格式搜索字体度量。这意味着 \LUATEX\ 查找 \type {OFMFONTS} 配置变量（如 \OMEGA\ 和 \ALEPH），而非 \type {TFMFONTS}（如 \TEX\ 和 \PDFTEX）。虚拟字体同理（\LUATEX\ 使用变量 \type {OVFFONTS} 而非 \type {VFFONTS}）。

\item 报告伸展或收缩次序的原语报告的值在 0 到 4 的便利范围内。由于某些宏包可能因此出错，我们还提供 \type {\eTeXgluestretchorder} 和 \type {\eTeXglueshrinkorder}，它们报告与 \ETEX\ 兼容的值。（新的）\texttt{fi} 值报告为 \type {-1}（因此在 \texttt{ifcase} 测试中使用时，该值会使程序进入 \texttt{else} 分支）。

\end{itemize}

\subsection{与 \PDFTEX\ 1.40 相比的变化}

\topicindex {\PDFTEX}

因为我们要生成 \PDF，最自然的起点就是广受欢迎的 \PDFTEX\ 程序。我们继承了稳定特性，丢弃了大部分实验性代码，并将一些功能提升为 \LUATEX\ 核心功能，这进而触发了原语的重命名。

出于兼容性原因，我们仍提及 \type {\pdf...} 命令，但 \LUATEX\ 有不同的后端接口。取代这些原语的是三个接口原语：\lpr {pdfextension}、\lpr {pdfvariable} 和 \lpr {pdffeedback}，它们接受关键字和可选的后续参数（下文我们仍使用 \tex {pdf} 前缀名称作为参考）。这样我们可以在需要时扩展功能，而无需改动核心引擎。前端与后端尽可能解耦。

\begin{itemize}\item （实验性的）snap 节点支持已被移除，因为在节点处理和属性之上构建此功能更为自然。随之消失的相关原语有：\orm {pdfsnaprefpoint}、\orm {pdfsnapy} 和 \orm {pdfsnapycomp}。

\item （实验性的）节点周围特殊间距支持也已被移除。随之消失的相关原语有：\orm {pdfadjustinterwordglue}、\orm {pdfprependkern} 和 \orm {pdfappendkern}，以及五个辅助原语 \orm {knbscode}、\orm {stbscode}、\orm {shbscode}、\orm {knbccode} 和 \orm {knaccode}。

\item 许多 ``\PDFTEX\ 原语已被移除，因为它们可以使用 \LUA\ 实现：\orm {pdfelapsedtime}、\orm {pdfescapehex}、\orm {pdfescapename}、\orm {pdfescapestring}、\orm {pdffiledump}、\orm {pdffilemoddate}、\orm {pdffilesize}、\orm {pdfforcepagebox}、\orm {pdflastmatch}、\orm {pdfmatch}、\orm {pdfmdfivesum}、\orm {pdfmovechars}、\orm {pdfoptionalwaysusepdfpagebox}、\orm {pdfoptionpdfinclusionerrorlevel}、\orm {pdfresettimer}、\orm {pdfshellescape}、\orm {pdfstrcmp} 和 \orm {pdfunescapehex}。

\item 版本相关原语 \orm {pdftexbanner}、\orm {pdftexversion} 和 \orm {pdftexrevision} 不再存在，因为与 \PDFTEX\ 开发不再有任何关系。

\item 实验性的 snapper 机制已被移除，因此原语 \orm {pdfignoreddimen}、\orm {pdffirstlineheight}、\orm {pdfeachlineheight}、\orm {pdfeachlinedepth} 和 \orm {pdflastlinedepth} 也一并消失。

\item 实验性原语 \lpr {primitive}、\lpr {ifprimitive}、\lpr {ifabsnum} 和 \lpr {ifabsdim} 被提升为核心原语。带 \type {\pdf*} 前缀的原始版本不可用。

\item 由于 \LUATEX\ 有不同的图像管理子系统，在此期间与祖先版本产生了更多偏离。我们不追随 \PDFTEX\ 的变化。

\item 提供了两个额外的令牌列表 \orm {pdfxformresources} 和 \orm {pdfxformattr}，作为 \orm {pdfxform} 关键字的替代方案。

\item 图像说明还支持 \type {visiblefilename}、\type {userpassword} 和 \type {ownerpassword}。密码选项仅与加密的 \PDF\ 文件相关。

\item 当前版本的 \LUATEX\ 不再将嵌入的 \PDF\ 文件中的字体替换和/或合并到外层 \PDF\ 文档的字体中。这种功能回退可能是暂时的，取决于重写后的字体后端的形态。

\item 原语 \orm {pdfpagewidth} 和 \orm {pdfpageheight} 已被移除，因为 \lpr {pagewidth} 和 \lpr {pageheight} 承担了该用途。

\item 原语 \orm {pdfnormaldeviate}、\orm {pdfuniformdeviate}、\orm {pdfsetrandomseed} 和 \orm {pdfrandomseed} 已被提升为不带 \type {pdf} 前缀的核心原语，因此原始命令不再被识别。

\item 原语 \lpr {ifincsname}、\lpr {expanded} 和 \lpr {quitvmode} 现在是核心原语。

\item 由于 hz 与突出机制已成为核心的一部分，相关原语 \lpr {lpcode}、\lpr {rpcode}、\lpr {efcode}、\lpr {leftmarginkern}、\lpr {rightmarginkern} 被提升为核心原语。两个命令 \lpr {protrudechars} 和 \lpr {adjustspacing} 取代了带 \type {\pdf} 前缀的原始版本。

\item hz 优化代码已部分重写，因此我们不再需要创建额外的字体实例。前端与后端已解耦，并生成更高效的（\PDF）代码。

\item 当 \lpr {adjustspacing} 的值为 2 时，hz 优化将应用于字形和字距。当值为 3 时，仅处理字形。小于 2 的值禁用此功能。值为 1 时，字体扩展在 \TEX\ 常规段落断行例程将段落拆分为行之后应用。在这种情况下，断行与标准 \TEX\ 行为一致（与 \PDFTEX\ 相同）。

\item \lpr {tagcode} 原语被提升为核心原语。

\item \lpr {letterspacefont} 功能现在是核心的一部分，但不会被修改（改进）。我们仅为遗留用途提供它。

\item \orm {pdfnoligatures} 原语现在是 \lpr {ignoreligaturesinfont}。

\item \orm {pdfcopyfont} 原语现在是 \lpr {copyfont}。

\item \orm {pdffontexpand} 原语现在是 \lpr {expandglyphsinfont}。

\item 由于位置跟踪在 \DVI\ 模式下也可用，\lpr {savepos}、\lpr {lastxpos} 和 \lpr {lastypos} 命令现在取代了带 \type {pdf} 前缀的原始版本。

\item 内省原语 \type {\pdflastximagecolordepth} 和 \type {\pdfximagebbox} 已被移除。可以使用外部应用程序确定这些属性，或使用内置的 \type {img} 库。

\item 初始化器 \orm {pdfoutput} 已被 \lpr {outputmode} 取代，\orm {pdfdraftmode} 现在是 \lpr {draftmode}。

\item 像素乘数尺寸 \orm {pdfpxdimen} 去掉了前缀，现在称为 \lpr {pxdimen}。

\item 增加了一个额外的 \orm {pdfimageaddfilename} 选项，可用于阻止将文件名写入 \PDF\ 文件。

\item 原语 \orm {pdftracingfonts} 现在是 \lpr {tracingfonts}，因为它与后端无关。

\item 实验性原语 \orm {pdfinsertht} 保留为 \lpr {insertht}。

\item 对于哪些元数据写入 \PDF\ 文件，有了更多控制。

\item 原语提升为核心原语以及字体与后端的分离意味着初始化命名空间 \type {pdftex} 已不复存在。

\end{itemize}

一项变更涉及所谓的 xforms 和 ximages。在 \PDFTEX\ 中它们以所谓的 whatsits 实现。但与其他 whatsit 不同，它们具有尺寸，在例如计算最优断行时需要考虑。在 \LUATEX\ 中，它们现在被提升为一种特殊类型的规则节点，这简化了需要这些尺寸的代码。

提升的另一个原因是这些概念很有用。后端可以提供在多处使用已渲染内容的能力，图像也很常见。如前文相关章节中已提及，我们现在有：

\begin{longtable}{l l}
\toprule \LUATEX & \PDFTEX \\
\midrule
\lpr {saveboxresource}             & \orm {pdfxform}           \\
\lpr {saveimageresource}           & \orm {pdfximage}          \\
\lpr {useboxresource}              & \orm {pdfrefxform}        \\
\lpr {useimageresource}            & \orm {pdfrefximage}       \\
\lpr {lastsavedboxresourceindex}   & \orm {pdflastxform}       \\
\lpr {lastsavedimageresourceindex} & \orm {pdflastximage}      \\
\lpr {lastsavedimageresourcepages} & \orm {pdflastximagepages} \\
\bottomrule
\end{longtable}

有一些 \lpr {pdffeedback} 功能与此相关，但这些是典型的后端特定功能。返回的索引应被视为 ``只是一个数字，虽然它仍具有与之前相同的含义（与对象相关），但你不应依赖这一点。

突出检测机制略有增强，以支持稍复杂的情况。当识别出突出字符时，会跳过一些节点：

\begin{itemize}
\item 零 glue 
\item 惩罚值 
\item 空 discretionary 
\item 常规零 kern 
\item 零尺寸规则 
\item 周围为零的数学节点 
\item dir 节点 
\item 空水平列表 
\item 局部 par 节点 
\item insert、mark 与 adjust 
\item 边界 
\item whatsit 
\end{itemize}

由于这可能还不够，你还可以使用突出边界节点来使下一个节点被忽略。当值为 1 或 3 时，在定位左边界条件的测试中会忽略下一个节点。当值为 2 或 3 时，在定位右边界条件时会忽略前一个节点（搜索从右向左进行）。这允许突出与例如移入页边距的内容结合使用：

\begin{typing}
\protrusionboundary1!\quad«Who needs protrusion?»
\end{typing}

\subsection{与 \ALEPH\ RC4 相比的变化}

\topicindex {\ALEPH}

因为我们需要正确的方向性排版，\ALEPH\ 的机制看起来最具吸引力。它们与 \OMEGA\ 提供的机制相当接近，因此接下来所述适用于这两个程序。

\begin{itemize}\item 扩展的 16 位数学原语（\orm {omathcode} 等）已被移除。

\item \OCP\ 处理已被完全移除，因此以下原语也被移除：\orm {ocp}、\orm {externalocp}、\orm {ocplist}、\orm {pushocplist}、\orm {popocplist}、\orm {clearocplists}、\orm {addbeforeocplist}、\orm {addafterocplist}、\orm {removebeforeocplist}、\orm {removeafterocplist} 和 \orm {ocptracelevel}。

\item \LUATEX\ 仅理解 \ALEPH\ 的 16 种方向说明符中的 4 种：\type {TLT}（拉丁文）、\type {TRT}（阿拉伯文）、\type {RTT}（中文/日文/韩文）、\type {LTL}（蒙古文）。所有其他方向说明符都会产生错误。除了关键字驱动模型外，我们还提供整数驱动模型。

\item \ALEPH\ 的输入翻译未实现，相关原语不可用：\orm {DefaultInputMode}、\orm {noDefaultInputMode}、\orm {noInputMode}、\orm {InputMode}、\orm {DefaultOutputMode}、\orm {noDefaultOutputMode}、\orm {noOutputMode}、\orm {OutputMode}、\orm {DefaultInputTranslation}、\orm {noDefaultInputTranslation}、\orm {noInputTranslation}、\orm {InputTranslation}、\orm {DefaultOutputTranslation}、\orm {noDefaultOutputTranslation}、\orm {noOutputTranslation} 和 \orm {OutputTranslation}。

\item 修复了若干错误，理清了令人困惑的实现细节。

\item 方向说明的扫描器现在允许在方向完全解析后有一个可选空格。

\item \type {\textasciicircum\textasciicircum{}} 记法已扩展：\type {\textasciicircum\textasciicircum{}\textasciicircum\textasciicircum{}} 后期望四个十六进制字符，\type {\textasciicircum\textasciicircum{}\textasciicircum\textasciicircum{}\textasciicircum\textasciicircum{}} 后必须给出六个十六进制字符。原始 \TEX\ 解释对 \type {\textasciicircum\textasciicircum{}} 情形仍然有效，但四个和六个变体不进行回溯，即当后面没有跟随正确数量的十六进制数字时会发出错误消息。因为 \type{\textasciicircum\textasciicircum{}^} 是正常的 \TEX\ 情形，我们也不支持奇数个的 \type {\textasciicircum\textasciicircum{}\textasciicircum\textasciicircum{}^}。

\item 方向更改命令 {\it 紧接其后} 的 glue 不是合法断点。

\item 若干需要从右到左感知的机制已得到改进。例如公式编号的放置。

\item 页面尺寸相关原语 \lpr {pagewidth} 和 \lpr {pageheight} 已被提升为核心原语。\prm {hoffset} 和 \prm {voffset} 原语已修复。

\item 原语 \type {\charwd}、\type {\charht}、\type {\chardp} 和 \type {\charit} 已被移除，因为我们有 \ETEX\ 变体 \type {\fontchar*}。

\item 两个尺寸寄存器 \lpr {pagerightoffset} 和 \lpr {pagebottomoffset} 现在是核心原语。

\item 方向相关原语 \lpr {pagedir}、\lpr {bodydir}、\lpr {pardir}、\lpr {textdir}、\lpr {mathdir} 和 \lpr {boxdir} 现在是核心原语。

\item 原语提升为核心原语以及移除所有其他原语意味着早期 \LUATEX\ 版本提供的初始化命名空间 \type {aleph} 已不复存在。

\end{itemize}

以上可以总结为：我们取其 32 位方面和大部分方向机制，将其合并到 \PDFTEX\ 代码库中作为进一步开发的起点。然后我们简化了方向性、修复了它并将其开放。

\subsection{来自各处的变化}

\type {\partokenname} 和 \type {\partokencontext} 原语取自发布在实现者邮件列表上的 \PDFTEX\ 变更文件。它们在 \PDFTEX\ 手册中有说明，并归类为 \ETEX\ 扩展。

\subsection{与标准 \WEBC 相比的变化}

\topicindex {\WEBC}

编译框架是 \WEBC，我们继续使用它，但去掉了 \PASCAL\ 到 \CCODE\ 的步骤。该框架还提供一些通用功能，处理从文件读取字节和在 \TDS\ 中定位文件。以下是我们的不同之处：

\begin{itemize}\item 不支持 mltex。

\item 不支持 enctex。

\item 以下编码相关的命令行开关被静默忽略，即使在非 \LUA\ 模式下也是如此：\type {-8bit}、\type {-translate-file}、\type {-mltex}、\type {-enc} 和 \type {-etex}。

\item \prm {openout} whatsit 不写入日志文件。

\item 部分所谓的 \WEBC\ 扩展在非 \KPSE\ 模式下难以设置，因为不读取 \type {texmf.cnf}：\type {shell-escape} 处于关闭状态（但由于 \LUA\ 的 \type {os.execute}，这不是问题），并且对 \type {openin} 和 \type {openout} 的偏执检查不会发生。然而，\LUA\ 脚本可以通过重载 \type {io.open} 等来自行完成这些检查。

\item ``E 选项不做任何有用的事。

\end{itemize}

\section{后端原语}\label{backendprimitives}

\subsection{更少的原语}

\topicindex {backend}
\topicindex {\PDF+backend}

在前一节中我们提到一些 \PDFTEX\ 原语被移除，另一些被提升为 \LUATEX\ 核心原语。这只是故事的一部分。为了在代码中分离后端特定原语，这些命令现在仅由少数几个取代。在传统 \TEX\ 中我们只有 \DVI\ 后端，但现在我们有两个：\DVI\ 和 \PDF。额外功能以 \TEX\ 术语中的 ``扩展 形式实现。通过更严格地分离，我们能够保持核心（前端）干净稳定，并隔离这些扩展。如果因某种原因需要额外的后端选项，可以在不触及核心的情况下实现。三个 \PDF\ 后端相关原语是：

\begin{typing}
\pdfextension command [specification]
\pdfvariable  name
\pdffeedback  name
\end{typing}

扩展会根据所给命令触发进一步解析。变量是一种寄存器，可以读写，而反馈则是报告某些信息（由于来自后端，它通常是一个令牌序列）。

\subsection{\lpr{pdfextension}\label{sec:pdfextensions}}

为了让 \LUATEX\ 超越单纯的 \TEX，你需要启用原语。从一开始就是如此。如果你想要传统的 \PDFTEX\ 原语（只要其功能仍然存在），现在可以这样做：

\begin{typing}
\def\pdfliteral             {\pdfextension literal}
\def\pdflateliteral         {\pdfextension lateliteral}
\def\pdfcolorstack          {\pdfextension colorstack}
\def\pdfsetmatrix           {\pdfextension setmatrix}
\def\pdfsave                {\pdfextension save\relax}
\def\pdfrestore             {\pdfextension restore\relax}
\def\pdfobj                 {\pdfextension obj }
\def\pdfrefobj              {\pdfextension refobj }
\def\pdfannot               {\pdfextension annot }
\def\pdfstartlink           {\pdfextension startlink }
\def\pdfendlink             {\pdfextension endlink\relax}
\def\pdfoutline             {\pdfextension outline }
\def\pdfdest                {\pdfextension dest }
\def\pdfthread              {\pdfextension thread }
\def\pdfstartthread         {\pdfextension startthread }
\def\pdfendthread           {\pdfextension endthread\relax}
\def\pdfinfo                {\pdfextension info }
\def\pdfcatalog             {\pdfextension catalog }
\def\pdfnames               {\pdfextension names }
\def\pdfincludechars        {\pdfextension includechars }
\def\pdffontattr            {\pdfextension fontattr }
\def\pdfmapfile             {\pdfextension mapfile }
\def\pdfmapline             {\pdfextension mapline }
\def\pdftrailer             {\pdfextension trailer }
\def\pdfglyphtounicode      {\pdfextension glyphtounicode }
\def\pdfrunninglinkoff      {\pdfextension linkstate 1 }
\def\pdfrunninglinkon       {\pdfextension linkstate 0 }
\end{typing}

内省原语可以定义为：

\begin{typing}
\def\pdftexversion    {\numexpr\pdffeedback version\relax}
\defpdftex           {\pdffeedback revision}
\def\pdflastlink      {\numexpr\pdffeedback lastlink\relax}
\def\pdfretval        {\numexpr\pdffeedback retval\relax}
\def\pdflastobj       {\numexpr\pdffeedback lastobj\relax}
\def\pdflastannot     {\numexpr\pdffeedback lastannot\relax}
\def\pdfxformname     {\numexpr\pdffeedback xformname\relax}
\def\pdfcreationdate          {\pdffeedback creationdate}
\def\pdffontname      {\numexpr\pdffeedback fontname\relax}
\def\pdffontobjnum    {\numexpr\pdffeedback fontobjnum\relax}
\def\pdffontsize      {\dimexpr\pdffeedback fontsize\relax}
\def\pdfpageref       {\numexpr\pdffeedback pageref\relax}
\def\pdfcolorstackinit        {\pdffeedback colorstackinit}
\end{typing}

配置相关的寄存器变为：

\begin{typing}
\edef\pdfcompresslevel            {\pdfvariable compresslevel}
\edef\pdfobjcompresslevel         {\pdfvariable objcompresslevel}
\edef\pdfrecompress               {\pdfvariable recompress}
\edef\pdfdecimaldigits            {\pdfvariable decimaldigits}
\edef\pdfgamma                    {\pdfvariable gamma}
\edef\pdfimageresolution          {\pdfvariable imageresolution}
\edef\pdfimageapplygamma          {\pdfvariable imageapplygamma}
\edef\pdfimagegamma               {\pdfvariable imagegamma}
\edef\pdfimagehicolor             {\pdfvariable imagehicolor}
\edef\pdfimageaddfilename         {\pdfvariable imageaddfilename}
\edef\pdfpkresolution             {\pdfvariable pkresolution}
\edef\pdfpkfixeddpi               {\pdfvariable pkfixeddpi}
\edef\pdfinclusioncopyfonts       {\pdfvariable inclusioncopyfonts}
\edef\pdfinclusionerrorlevel      {\pdfvariable inclusionerrorlevel}
\edef\pdfignoreunknownimages      {\pdfvariable ignoreunknownimages}
\edef\pdfgentounicode             {\pdfvariable gentounicode}
\edef\pdfomitcidset               {\pdfvariable omitcidset}
\edef\pdfomitcharset              {\pdfvariable omitcharset}
\edef\pdfomitinfodict             {\pdfvariable omitinfodict}
\edef\pdfomitinfodict             {\pdfvariable omitinfodict}
\edef\pdfomitmediabox             {\pdfvariable omitmediabox}
\edef\pdfomitprocset              {\pdfvariable omitprocset}
\edef\pdfptexprefix               {\pdfvariable ptexprefix}
\edef\pdfpagebox                  {\pdfvariable pagebox}
\edef\pdfminorversion             {\pdfvariable minorversion}
\edef\pdfuniqueresname            {\pdfvariable uniqueresname}

\edef\pdfhorigin                  {\pdfvariable horigin}
\edef\pdfvorigin                  {\pdfvariable vorigin}
\edef\pdflinkmargin               {\pdfvariable linkmargin}
\edef\pdfdestmargin               {\pdfvariable destmargin}
\edef\pdfthreadmargin             {\pdfvariable threadmargin}
\edef\pdfxformmargin              {\pdfvariable xformmargin}

\edef\pdfpagesattr                {\pdfvariable pagesattr}
\edef\pdfpageattr                 {\pdfvariable pageattr}
\edef\pdfpageresources            {\pdfvariable pageresources}
\edef\pdfxformattr                {\pdfvariable xformattr}
\edef\pdfxformresources           {\pdfvariable xformresources}
\edef\pdfpkmode                   {\pdfvariable pkmode}

\edef\pdfsuppressoptionalinfo     {\pdfvariable suppressoptionalinfo }
\edef\pdftrailerid                {\pdfvariable trailerid }
\end{typing}

这些变量是内部变量，因此是匿名的。当你询问先前定义的几个变量的含义时：

\begin{typing}
\meaning\pdfhorigin
\meaning\pdfcompresslevel
\meaning\pdfpageattr
\end{typing}

你会得到：

\begin{typing}
macro:->[internal backend dimension]
macro:->[internal backend integer]
macro:->[internal backend tokenlist]
\end{typing}

\prm {edef} 也可以是 \prm {def}，但预先展开查找相关寄存器效率更高。

后端派生自 \PDFTEX，因此适用相同的语法。不过，\type {outline} 命令接受一个 \type {objnum} 后跟一个数字。不进行检查，因此使用时最好是一个有效的（已刷新的）对象。

为了（或多或少）与 \PDFTEX\ 兼容，我们也支持抑制某些信息的选项，但通过位集实现：

\begin{typing}
\pdfvariable suppressoptionalinfo \numexpr
        0
    +   1   % PTEX.FullBanner
    +   2   % PTEX.FileName
    +   4   % PTEX.PageNumber
    +   8   % PTEX.InfoDict
    +  16   % Creator
    +  32   % CreationDate
    +  64   % ModDate
    + 128   % Producer
    + 256   % Trapped
    + 512   % ID
\relax
\end{typing}

此外，你可以重载 trailer id，但我们不做任何有效性检查，因此你必须传入一个有效的数组。以下类似引擎通常生成的形式。这里你甚至需要包含括号！

\begin{typing}
\pdfvariable trailerid {[
    <FA052949448907805BA83C1E78896398>
    <FA052949448907805BA83C1E78896398>
]}
\end{typing}

虽然我们从合并 \PDFTEX\ 和 \ALEPH\ 起步，但到现在为止代码库和功能已偏离了这些来源。这里我们展示可以传递给扩展的选项。\type {shipout} 选项是一个兼容性特性。相反，可以使用 \type {deferred} 前缀。

\begin{typing}
\pdfextension literal
    [shipout] [ direct | page | raw ] { tokens }
\end{typing}

\begin{typing}
\pdfextension dest
    num integer | name { tokens }!crlf
    [ fitbh | fitbv | fitb | fith| fitv | fit |
      fitr <rule spec> | xyz [ zoom <integer> ]
\end{typing}

\begin{typing}
\pdfextension annot
    reserveobjnum | useobjnum <integer>
    { tokens }
\end{typing}

\begin{typing}
\pdfextension save
\end{typing}

\begin{typing}
\pdfextension restore
\end{typing}

\begin{typing}
\pdfextension setmatrix
    { tokens }
\end{typing}

\begin{typing}
[ \immediate ] \pdfextension obj
    reserveobjnum
\end{typing}

\begin{typing}
[ \immediate ] \pdfextension obj
    [ useobjnum <integer> ]
    [ uncompressed ]
    [ stream  [ attr { tokens } ] ]
    [ file ]
    { tokens }
\end{typing}

\begin{typing}
\pdfextension refobj
    <integer>
\end{typing}

\begin{typing}
\pdfextension colorstack
    <integer>
    set { tokens } | push { tokens } | pop | current
\end{typing}

\begin{typing}
\pdfextension startlink
    [ attr { tokens } ]
    user { tokens } | goto | thread
    [ file { tokens } ]
    [ page <integer> { tokens } | name { tokens } | num  integer ]
    [ newwindow  | nonewwindow ]
\end{typing}

\begin{typing}
\pdfextension endlink
\end{typing}

\begin{typing}
\pdfextension startthread
    num <integer> | name { tokens }
\end{typing}

\begin{typing}
\pdfextension endthread
\end{typing}

\begin{typing}
\pdfextension thread
    num <integer> | name { tokens }
\end{typing}

\begin{typing}
\pdfextension outline
    [ attr { tokens } ]
    [ useobjnum <integer> ]
    [ count <integer> ]
    { tokens }
\end{typing}

\begin{typing}
\pdfextension glyphtounicode
    { tokens }
    { tokens }
\end{typing}

\begin{typing}
\pdfextension catalog
    { tokens }
    [ openaction
      user { tokens } | goto | thread
      [ file { tokens } ]
      [ page <integer> { tokens } | name { tokens } | num <integer> ]
      [ newwindow  | nonewwindow ] ]
\end{typing}

\begin{typing}
\pdfextension fontattr
    <integer>
    {tokens}
\end{typing}

\begin{typing}
\pdfextension mapfile
    {tokens}
\end{typing}

\begin{typing}
\pdfextension mapline
    {tokens}
\end{typing}

\begin{typing}
\pdfextension includechars
    {tokens}
\end{typing}

\begin{typing}
\pdfextension info
    {tokens}
\end{typing}

\begin{typing}
\pdfextension names
    {tokens}
\end{typing}

\begin{typing}
\pdfextension trailer
    {tokens}
\end{typing}

\subsection{默认值}

引擎设置以下默认值。

\begin{typing}
\pdfcompresslevel         9
\pdfobjcompresslevel      1 % used: (0,9)
\pdfrecompress            0 % mostly for debugging
\pdfdecimaldigits         4 % used: (3,6)
\pdfgamma              1000
\pdfimageresolution      71
\pdfimageapplygamma       0
\pdfimagegamma         2200
\pdfimagehicolor          1
\pdfimageaddfilename      1
\pdfpkresolution         72
\pdfpkfixeddpi            0
\pdfinclusioncopyfonts    0
\pdfinclusionerrorlevel   0
\pdfignoreunknownimages   0
\pdfgentounicode          0
\pdfomitcidset            0
\pdfomitcharset           0
\pdfomitinfodict          0
\pdfomitmediabox          0
\pdfomitprocset           0
\pdfptexprefix            0
\pdfpagebox               0
\pdfminorversion          4
\pdfuniqueresname         0

\pdfhorigin             1in
\pdfvorigin             1in
\pdflinkmargin          0pt
\pdfdestmargin          0pt
\pdfthreadmargin        0pt
\pdfxformmargin         0pt
\end{typing}

\subsection{向后兼容}

如果你还想要一些向后兼容，可以添加：

\begin{typing}
\let\pdfpagewidth      \pagewidth
\let\pdfpageheight     \pageheight

\let\pdfadjustspacing  \adjustspacing
\let\pdfprotrudechars  \protrudechars
\let\pdfnoligatures    \ignoreligaturesinfont
\let\pdffontexpand     \expandglyphsinfont
\let\pdfcopyfont       \copyfont

\let\pdfxform          \saveboxresource
\let\pdflastxform      \lastsavedboxresourceindex
\let\pdfrefxform       \useboxresource

\let\pdfximage         \saveimageresource
\let\pdflastximage     \lastsavedimageresourceindex
\let\pdflastximagepages\lastsavedimageresourcepages
\let\pdfrefximage      \useimageresource

\let\pdfsavepos        \savepos
\let\pdflastxpos       \lastxpos
\let\pdflastypos       \lastypos

\let\pdfoutput         \outputmode
\let\pdfdraftmode      \draftmode

\let\pdfpxdimen        \pxdimen

\let\pdfinsertht       \insertht

\let\pdfnormaldeviate  \normaldeviate
\let\pdfuniformdeviate \uniformdeviate
\let\pdfsetrandomseed  \setrandomseed
\let\pdfrandomseed     \randomseed

\let\pdfprimitive      \primitive
\let\ifpdfprimitive    \ifprimitive

\let\ifpdfabsnum       \ifabsnum
\let\ifpdfabsdim       \ifabsdim
\end{typing}

甚至：

\begin{typing}
\newdimen\pdfeachlineheight
\newdimen\pdfeachlinedepth
\newdimen\pdflastlinedepth
\newdimen\pdffirstlineheight
\newdimen\pdfignoreddimen
\end{typing}

\section{方向}

\topicindex {\OMEGA}
\topicindex {\ALEPH}
\topicindex {directions}

\subsection{四种方向}

\LUATEX\ 中的方向模型继承自 \OMEGA/\ALEPH，但我们尝试做了一些改进。曾有一段时间我们尝试恢复模式，但后来发现它与嵌套方向冲突而禁用了它。这本身的副作用是节点列表在方向节点方面不再平衡，当一系列 dir 更改在没有分组的情况下发生时，这可能产生副作用。

在扩展 \PDF\ 后端以支持方向时发现了一些不一致之处，因此我们决定仅支持四种有意义的模型：\type {TLT}（拉丁文）、\type {TRT}（阿拉伯文）、\type {RTT}（中文/日文/韩文）和 \type {LTL}（蒙古文）。

\subsection{工作原理}

我们的方法是再次使列表平衡，但尽量避免一些副作用。如果忽略空格（转换为 glue），发生的情况相当直观，但即使在那里，仔细看也能理解。然而，当没有使用适当的嵌套分组时，这种逻辑使组内切换几乎毫无用处：从右到左多次嵌套切换会导致间距因嵌套镜像而一个接一个地排列。当然，一个合理的宏包会为用户管理这些，但我们这里讨论的是底层 dir 注入。

情况如下：

\begin{typing}
textdir TRT nur {textdir TLT run textdir TRT NUR} nur
\end{typing}

逐步变为：

\begin{typing}
injected: [+TRT]nur {[+TLT]run [+TRT]NUR} nur
balanced: [+TRT]nur {[+TLT]run [-TLT][+TRT]NUR[-TRT]} nur[-TRT]
result  : run {RUNrun } run
\end{typing}

而这个：

\begin{typing}
textdir TRT nur {nur textdir TLT run textdir TRT NUR} nur
\end{typing}

变为：

\begin{typing}
injected: [+TRT]nur {nur [+TLT]run [+TRT]NUR} nur
balanced: [+TRT]nur {nur [+TLT]run [-TLT][+TRT]NUR[-TRT]} nur[-TRT]
result  : run {run RUNrun } run
\end{typing}

现在，在以下示例中注意我们花括号的位置：

\begin{typing}
textdir TRT nur {{textdir TLT run} {textdir TRT NUR}} nur
\end{typing}

这变为：

\begin{typing}
[def]
\def\ltr{\linedir TLT\relax}
\def\rtl{\linedir TRT\relax}
\end{typing}



而这一次：

\begin{typing}
run {\rtl nur {\ltr run \rtl NUR \ltr run \rtl NUR} nur}
run {\ltr run {\rtl nur \ltr RUN \rtl nur \ltr RUN} run}
\end{typing}



输出为间距正确的：

任何比这更复杂的情况，如 skip 和 penalty 的组合或 kern，都应在输入或宏包中处理，因为我们无法预测预期行为。事实上，\lpr {linedir} 只是一个便利的附加功能，也可以使用节点列表解析来实现。

方向的复杂性在于它们通常需要跨组工作，因此使用了单独的分组相关堆栈。副作用是可能存在仅包含一个局部 par 节点、后跟方向同步节点的段落。这样的段落被视为空段落，因此被忽略。因为 \type {\noindent} 不注入任何内容，而 \type {\indent} 会注入一个 box，所以仅含缩进和方向的段落被视为有内容的段落处理。

默认情况下，包含 dir 节点的显示公式前的段落从不被忽略。更改此行为可能会破坏现有文档，但当你将 \lpr {mathemptydisplaymode} 设为 \type {1} 时，显示公式前的空段落将被忽略。

\subsection{用 \lpr {breakafterdirmode}}

dir 节点后的 glue 在断行决策中被忽略，但你可以通过将 \lpr {breakafterdirmode} 设为 \type {1} 来绕过。下表对比了不同设置下文本方向切换处的断行行为：

\begin{longtable}{p{5cm} p{5cm}}
\toprule
\textbf{breakafterdirmode=0} & \textbf{breakafterdirmode=1} \\
\midrule
pre \{textdir TLT xxx\} post & pre \{textdir TLT xxx\} post \\
pre \{textdir TLT xxx \}post & pre \{textdir TLT xxx \}post \\
pre\{ textdir TLT xxx\} post & pre\{ textdir TLT xxx\} post \\
pre\{ textdir TLT xxx \}post & pre\{ textdir TLT xxx \}post \\
pre \{ textdir TLT xxx \} post & pre \{ textdir TLT xxx \} post \\
\bottomrule
\end{longtable}

\subsection{用 \lpr {shapemode}}

对 \ALEPH\ 方向模型的另一项调整是对由 \prm {hangindent} 和 \prm {parshape} 驱动的形状进行控制。这由一个新参数 \lpr {shapemode} 控制：

\begin{longtable}{c l l}
\toprule 值    & \prm {hangindent} & \prm {parshape} \\
\midrule
\type{0} &  正常             & 正常            \\
\type{1} &  镜像           & 正常            \\
\type{2} &  正常             & 镜像          \\
\type{3} &  镜像           & 镜像          \\
\bottomrule
\end{longtable}

该值在段落结束后重置为零（如 \prm {hangindent} 和 \prm {parshape}）。你可以使用负值来防止这种情况。\ref{fig:shapemode} 中给出了几个示例。

\begin{figure}[htbp]
\centering
\fbox{\parbox{3.5cm}{\centering TLT 方向\\hangindent=40pt\\hangafter=$-3$}}
\quad
\fbox{\parbox{3.5cm}{\centering TLT 方向\\parshape 4 行\\不规则缩进}}
\quad
\fbox{\parbox{3.5cm}{\centering TRT 方向\\hangindent=40pt\\shapemode=0}}
\\[6pt]
\fbox{\parbox{3.5cm}{\centering TRT 方向\\parshape 4 行\\shapemode=0}}
\quad
\fbox{\parbox{3.5cm}{\centering TRT 方向\\hangindent=40pt\\shapemode=3}}
\quad
\fbox{\parbox{3.5cm}{\centering TRT 方向\\parshape 4 行\\shapemode=3}}
\caption{\type{shapemode} 的效果。}\label{fig:shapemode}
\end{figure}

\subsection{符号或数字}

内部实现与 \ALEPH\ 不同。首先，我们不使用 whatsit 而是专用节点，而且我们只有 4 个方向映射到 4 个数字。文本方向节点可以标记节点序列的开始或结束，因此有两种状态。在 \TEX\ 端我们看不到这些状态，因为 \TEX\ 本身会在需要时添加适当的结束状态节点。

符号名 \type {TLT}、\type {TRT} 等源自 \OMEGA。在 \LUATEX\ 中我们还有基于数字的模型，有时更有意义。

\begin{longtable}{c l l}
\toprule 值     & 等价 \\
\midrule
\type {0} & TLT \\
\type {1} & TRT \\
\type {2} & LTL \\
\type {3} & RTT \\
\bottomrule
\end{longtable}

我们支持 \OMEGA\ 原语 \orm {textdir}、\orm {pardir}、\orm {pagedir}、\orm {pardir} 和 \orm {mathdir}。它们接受三字符关键字。通过数字设置方向的原语有：\lpr {textdirection}、\lpr {pardirection}、\lpr {pagedirection}、\lpr {bodydirection} 和 \lpr {mathdirection}。为 box 指定方向时，可以使用 \type {bdir} 代替 \type {dir}。

\section{实现说明}

\subsection{内存分配}

\topicindex {memory}

传统 \TEX\ 用于令牌和节点的单一内部内存堆被拆分为两个独立数组。每个数组在需要时动态增长。

与主内存相关的 \type {texmf.cnf} 设置不再使用（这些是：\type {main_memory}、\type {mem_bot}、\type {extra_mem_top} 和 \type {extra_mem_bot}）。``主内存不足 错误仍可能发生，但限制因素现在是系统中的 RAM 容量，而非预定义限制。

此外，节点的内存（解）分配例程已完全重写。相关代码现在位于 C 文件 \type {texnode.c} 中，基本上使用十几个 ``avail 列表而非双链接模型。增加了一个额外的函数层，使代码可以按类型请求节点，而不是直接申请一定数量的内存字。

由于拆分为两个数组以及由此产生的数据结构差异，部分宏已被复制。例如，现在有 \type {vlink} 和 \type {vinfo}，以及 \type {token_link} 和 \type {token_info}。对变量内存数组的所有访问现在都隐藏在一个名为 \type {vmem} 的宏之后。我们提到这一点是因为使用 \TEX book 作为参考仍然相当有效，但不适用于内存相关细节。另一个重要细节是我们有双链接节点列表，且大多数节点携带更多数据。

输入行缓冲区和池大小现在也在需要时重新分配，\type {texmf.cnf} 设置 \type {buf_size} 和 \type {pool_size} 被静默忽略。

\subsection{稀疏数组}

\prm {mathcode}、\prm {delcode}、\prm {catcode}、\prm {sfcode}、\prm {lccode} 和 \prm {uccode}（以及新的 \lpr {hjcode}）表现在是用 \CCODE\ 实现的稀疏数组。它们不再是 \TEX\ ``等价表 的一部分，由于每个表有 110 万条目、每条目占几个内存字，这在内存使用上产生了重大差异。性能并未因此受到实质影响。

\prm {catcode}、\prm {sfcode}、\prm {lccode}、\prm {uccode} 和 \lpr {hjcode} 赋值在使用 \ETEX\ 跟踪例程 \prm {tracingassigns} 和 \prm {tracingrestores} 时不显示，但我们不认为这是真正的限制。

当前实现的一个副作用是 \prm {global} 现在在处理上比非全局赋值更昂贵，但不会有多少用户注意到。

字体内的 glyph id 也通过稀疏数组管理，因为 glyph id 可高达索引 $2^{21}-1$，但从不直接访问，因此用户同样不会注意到。

\subsection{简单的单字符 csname}

\topicindex {csnames}

单字符命令在内部不再被特殊对待，它们像多字符 csname 一样存储在哈希中。

显示控制序列的代码在决定是否添加尾随空格时，会显式检查长度是否为 1。

活动字符在内部实现为一种特殊类型的多字符控制序列，使用一个在其他情况下不可能获得的前缀。

\subsection{压缩的格式文件}

\topicindex {format}

格式通过 \type {zlib} 处理，使其可缩小到未压缩时约一半的大小。这会消耗稍多的 \CPU\ 周期，但磁盘 I/O\ 少得多，因此总体应该更快。我们使用 3 级压缩，我们发现这是文件大小与解压速度之间的最佳权衡。

\subsection{二进制文件读取}

\topicindex {files+binary}

所有内部代码都做了修改：如果 \type {read_xxx_file} 回调之一未设置，则文件由 \CCODE\ 函数读取，基本使用与回调相同的约定：一次性读入足够容纳整个文件内容的缓冲区。虽然这比以前的代码（主要使用 \type {getc} 调用）使用更多内存，但可能快不少（取决于你的 I/O\ 子系统）。

\subsection{制表符与空格}

\topicindex {space}
\topicindex {newline}

我们遵循其他 \TEX\ 引擎处理尾随制表符和空格的方式。几十年来，尾随制表符和空格（换行符之前）会从输入中移除，但此行为在 2017 年 9 月被修改为仅处理空格。我们知道这可能在现有工作流中引入兼容性问题，但因为我们不想与上游 \TEXLIVE\ 有太多差异，所以直接跟进了该补丁（这是功能性变更而非真正的修复）。处理可能的兼容性问题取决于宏包维护者，在 \LUATEX\ 中他们可以通过处理文件读取的回调来做到。

之前的行为是已知的副作用，而且（由于这类输入通常来自生成的源），如果空格和/或制表符是有意保留的，通常通过在该行添加注释令牌来处理。我们知道这与我们的其他一些选择相矛盾，但与其他引擎保持一致以及在 \KPSE\ 模式下使用公共文件 I/O\ 层这一事实，可能会产生破坏兼容性的副作用。我们仍然坚持在日志级别上我们可以（也可能）更不兼容的观点。我们已经暴露了更多细节。

\subsection{超链接}

\topicindex {hyperlinks}

有一个实验性功能使多行超链接表现稍好，修复了在从右到左排版中出现的一些副作用，但这些副作用也可能出现在从左到右排版中。由于此问题直到 2023 年才被注意到，且它在一定程度上取决于宏包如何处理超链接，因此该修复目前由参数控制：

\begin{typing}
\pdfvariable linking = 1
\end{typing}

这样（我们希望）遗留文档能按预期输出，无论预期是什么。处理的方面之一涉及（不寻常的）左右 skip。
